\chapter{Testing and Evaluation}

\section{Testing Strategy}

The evaluation approach is organized around business risk rather than around
implementation layers alone. The most important questions are whether the
gateway preserves correct transaction state, protects access boundaries, keeps
merchant data scoped correctly, and records enough evidence for operational
review.

\begin{longtable}{L{0.22\textwidth}L{0.36\textwidth}L{0.32\textwidth}}
\caption{Testing and evaluation matrix}\\
\toprule
\textbf{Category} & \textbf{Coverage} & \textbf{Evidence type} \\
\midrule
\endfirsthead
\toprule
\textbf{Category} & \textbf{Coverage} & \textbf{Evidence type} \\
\midrule
\endhead
Business rule validation & Merchant readiness, QR account readiness, VietQR generation, duplicate handling, refund eligibility, expiration, notification retry, reconciliation, and audit-aware state changes. & Automated business rule checks. \\
\bodyrowrule
End-to-end scenario validation & Merchant onboarding, QR account configuration, payment creation, callback processing, refund finalization, notification delivery, late evidence handling, and manual recovery. & Automated end-to-end transaction scenarios. \\
\bodyrowrule
Internal operations validation & Internal sign-in, merchant lifecycle actions, QR account management, review workflows, reconciliation handling, and evidence exploration. & Automated operational workflow checks and manual review walkthroughs. \\
\bodyrowrule
Merchant dashboard validation & Merchant sign-in, analytics windows, scoped list and detail views, QR image and payload display, profile visibility, credential metadata visibility, and password maintenance. & Automated dashboard interaction tests and scoped data checks. \\
\bodyrowrule
Manual usability review & Readability, layout stability, chart clarity, mobile behavior, and operator usability. & Manual browser review. \\
\bottomrule
\end{longtable}

\section{Business Rule Evaluation}

The strongest part of the evaluation is the coverage of transaction rules and
state transitions. Payment creation, refund eligibility, callback handling,
notification retry, and reconciliation behavior are all checked through
automated scenarios that focus on expected business outcomes rather than only on
technical success responses.

This is important because the system is not merely a request router. It must
decide when a merchant is ready, whether a merchant has an active VietQR
receiving account, when a duplicate payment may be reused, when a refund is
allowed, when a delayed callback becomes suspicious, and when a failed
notification is safe to retry.

The VietQR pilot adds tests for QR-specific failure and success paths:

\begin{itemize}
  \item payment creation fails when an active merchant has no active QR account;
  \item VND and whole-amount constraints are enforced before QR generation;
  \item generated payments include QR reference, VietQR payload, and PNG data
  URL evidence;
  \item duplicate equivalent pending payment requests return the same stored QR
  evidence;
  \item provider callbacks reject missing, invalid, or stale HMAC signatures.
\end{itemize}

\section{Dashboard And Access Evaluation}

The report also evaluates whether the two human-facing dashboards enforce the
correct boundaries. Internal users must be able to review and recover
operational issues, while merchant users must only see records that belong to
their own merchant account.

Dashboard evaluation therefore focuses on:

\begin{itemize}
  \item sign-in and session behavior for internal and merchant-facing users;
  \item merchant-scoped visibility for analytics and record exploration;
  \item visibility of useful metadata without exposure of raw secret material;
  \item internal-only QR account controls and merchant-facing QR evidence
  display;
  \item practical usability for review, investigation, and account maintenance
  tasks.
\end{itemize}

\section{Pilot Demo And Smoke Evaluation}

The branch includes a visual demo merchant, real-HTTP smoke scripts, and the
existing Postman/Newman scenario. Together they exercise the same flow used in
demonstration:

\begin{enumerate}
  \item internal login, merchant onboarding, credential creation, and QR account
  setup;
  \item merchant HMAC payment creation returning \code{qr_reference},
  \code{qr_content}, and \code{qr_image_base64};
  \item signed provider callback that finalizes the payment;
  \item refund creation and signed refund callback processing;
  \item worker-driven webhook delivery or retry behavior;
  \item merchant-side webhook verification, duplicate event handling, and
  visible checkout transition to success or failure.
\end{enumerate}

The commands \code{scripts/smoke_e2e_demo.py --outcome success} and
\code{scripts/smoke_e2e_demo.py --outcome failed} provision unique merchants
through the real Ops API, create QR payments through merchant HMAC, call the
provider callback surface, and poll until the worker webhook reaches the demo
merchant. They also verify that the Merchant Dashboard API exposes the same
payment. The simulator response is explicitly checked while the demo order is
still pending, which proves that only the merchant webhook controls the final
checkout result.

This gives the demonstration a visible real-world story: the merchant backend
receives a QR image for checkout, the customer can scan it in a banking app,
the simulator replaces provider confirmation, and the gateway then records
evidence and notifies the merchant. Figure~\ref{fig:demo-merchant-checkout}
shows both the pending and webhook-confirmed states used in browser review.

\section{Manual Review And Overall Assessment}

Manual browser review complements the automated evaluation by checking visual
and operational qualities that are difficult to reduce to pure rule-based
tests. These checks include page readability, table overflow, chart clarity,
mobile behavior, and general workflow stability.

Overall, the evaluation indicates that the Mini Payment Gateway covers its pilot
goals well:

\begin{itemize}
  \item money-flow rules are protected by explicit validation and state control;
  \item QR payment generation is tied to audited Ops configuration rather than
  uncontrolled merchant-side data;
  \item external evidence and merchant notification behavior are preserved as
  durable reviewable events;
  \item internal and merchant-facing access boundaries are treated as separate
  security concerns;
  \item both dashboards support the operational and merchant visibility goals of
  the project.
  \item the separate demo merchant demonstrates server-side HMAC integration
  and webhook-driven customer status without weakening the gateway's API-only
  boundary.
\end{itemize}

The main remaining limitation is typical for an academic pilot system: some
usability and layout quality still depend on manual review in addition to
automated validation.
